% Pillar-5 (MIN-REPORT-RL) reference verifier: from checklist to executable CI gate
\section{A Reference Verifier: MIN-REPORT-RL as an Executable Gate}
\label{sec:p5-verifier}
The toolchain of \secref{sec:p5-toolchain} is a specification; a standard is
only enforceable once a machine can reject a non-compliant run without a human
in the loop. We therefore implement a reference verifier
(\texttt{registry/provenance/minreport.py}; signature checks use the
\texttt{cryptography} package) that turns the
seven-field run-manifest block from a descriptive checklist into an automated gate. This is
the distinction from the P6 catalog: a registry \emph{catalogs} what a run declares; the verifier
\emph{recomputes} and \emph{rejects}.

\paragraph{The record.} A run emits a provenance record
(\texttt{*.provenance.json}, schema \texttt{MIN-REPORT-RL/2.0}) whose blocks
include \texttt{min\_report\_rl} (the seven run-manifest items---loss form,
reference-KL, sampler backend, telemetry, group-size schedule, held-out split,
decontamination); \texttt{provenance} (anchors: \texttt{git\_sha},
\texttt{code\_file} with its \texttt{code\_hash}, \texttt{reward\_fn\_hash},
\texttt{config\_hash}); \texttt{rigor} (\texttt{n\_seeds} and the seed list);
\texttt{evidence} (repository-local, SHA-256-pinned run, split,
decontamination, and analysis artifacts); and \texttt{results}. A
\texttt{generate} command drafts the record from the run's configuration; the
draft is unsigned until it passes the checks below.

\paragraph{What the verifier checks.} \texttt{minreport.py verify} runs a
fixed list of checks and reports a reporting-completeness percentage
separately from acceptance. \emph{Schema and completeness}: the record
validates against the versioned schema and all seven run-manifest items are
present. \emph{Integrity}: the recorded \texttt{code\_hash},
\texttt{reward\_fn\_hash}, and \texttt{config\_hash} are \emph{recomputed}
against the live repository, and \texttt{git\_sha} must be the full commit
of the checked-out tree, catching drift between the run and the claim.
\emph{Rigor}: at least three distinct integer seeds, one completed run
artifact per seed, a saved held-out split with its count, and saved analysis
inputs with uncertainty. A record is \textsc{accepted} (grade A) only if every
check passes \emph{and} it carries an Ed25519 signature under a key the caller
supplies; otherwise it is \textsc{rejected}, and the command exits non-zero so
a continuous-integration job can block a merge. Three seeds are a replication
rule, not a power claim, and a signature establishes approval by a key, not
that execution occurred.

\paragraph{Applied to this repository's own records.} We re-ran the current
verifier in diagnostic mode on the two provenance records in the repository
(\texttt{registry/provenance/\{campaign,sweep\}.provenance.json};
\texttt{tools/review\_ledger/P5\_minreport\_verifier\_demo.py}). Both are
\textsc{rejected} at $16\%$ reporting completeness: they predate the 2.0
schema, carry an abbreviated \texttt{git\_sha}, and pin no evidence
artifacts. The seed rule separates them as intended---the multi-seed campaign
passes the three-seed check and the single-seed sweep fails it---but neither
record is publication-ready. An earlier revision reported grade~A for the
campaign record under a previous verifier; that result does not reproduce and
is withdrawn.

\paragraph{Tamper-evidence.} A companion signer
(\texttt{registry/provenance/sign.py}) produces a detached Ed25519 signature
over the canonicalized record (sorted keys, ASCII JSON) with no algorithm
fallback; its CLI refuses to sign a record that fails the evidence checks. In
our re-run, a signature made over the campaign record verifies against the
pristine copy and fails after a single field (\texttt{git\_sha}) is altered.
This is the same detached-attestation idea that software supply-chain
frameworks such as in-toto (Torres-Arias et al., USENIX Security 2019) apply to
build steps; here it covers the run record only, not the execution that
produced it.
